\subsection{域上代数的张量积}

设 K 是一个交换域，E 和 F 是 K 上的两个代数，其各自的单位元 \(e, e'\) 都非零。则同态 \(\eta_{\mathbf{E}}: \mathbf{K} \to \mathbf{E}\) 和 \(\eta_{\mathbf{F}}: \mathbf{K} \to \mathbf{F}\) （§1, 第 3 号）是单射，这使我们可以把 K 等同于 E（相应地，F）的一个子域。典范同态 \(u: \mathbf{E} \to \mathbf{E} \otimes_{\mathbf{K}} \mathbf{F}\) 和 \(v: \mathbf{F} \to \mathbf{E} \otimes_{\mathbf{K}} \mathbf{F}\) （由 \(u(x) = x \otimes e'\) 和 \(v(y) = e \otimes y\) 定义）是单射（II, §7, 第 9 号, 命题 19），这使我们可以把 E 和 F 等同为 \(\mathbf{E} \otimes_{\mathbf{K}} \mathbf{F}\) 的子代数，两者的单位元都是单位元 \(e \otimes e'\) （ \(\mathbf{E} \otimes_{\mathbf{K}} \mathbf{F}\) 的）。在 \(\mathbf{E} \otimes_{\mathbf{K}} \mathbf{F}\) 中， \(\mathbf{E} \cap \mathbf{F} = \mathbf{K}\) （II, §7, 第 9 号, 命题 19）。

若 \(\mathbf{E}'\) 和 \(\mathbf{F}'\) 分别是 \(\mathbf{E}\) 和 \(\mathbf{F}\) 的子代数，则典范同态 \(\mathbf{E}' \otimes_{\mathbb{K}} \mathbf{F}' \to \mathbf{E} \otimes_{\mathbb{K}} \mathbf{F}\) 是单射，这使我们可以把 \(\mathbf{E}' \otimes_{\mathbb{K}} \mathbf{F}'\) 等同于 \(\mathbf{E} \otimes_{\mathbb{K}} \mathbf{F}\) 的由 \(\mathbf{E}' \cup \mathbf{F}'\) 生成的子代数（II, §7, 第 7 号, 命题 14）。

命题 6. 设 E, F 是交换域上的两个非零代数，C（相应地 D）是 E（相应地 F）的一个子代数，且 C'（相应地 D'）是 C 在 E 中（相应地 D 在 F 中）的中心化子。则 C ⊗\_K D 在 E ⊗\_K F 中的中心化子是 C' ⊗\_K D'。

这归结为验证：元素 \(z = \sum_{i} x_{i} \otimes y_{i}\) （属于 \(C \otimes_{K} D\) 的中心化子， \(x_{i} \in F, y_{i} \in F\) ）含于 \(C' \otimes_{K} D'\) ；我们知道

\[
\mathrm{C} ^ {\prime} \otimes_ {\mathbf {K}} \mathrm{D} ^ {\prime} = (\mathrm{C} ^ {\prime} \otimes_ {\mathbf {K}} \mathrm{F}) \cap (\mathrm{E} \otimes_ {\mathbf {K}} \mathrm{D} ^ {\prime})
\]

（II, §7, 第 7 号, 命题 14 的推论）。可假定诸 \(y_{i}\) 在 \(\mathbf{K}\) 上线性无关；对所有 \(x \in \mathbb{C}\) ，必有 \((x \otimes e')z = z(x \otimes e')\) ，即 \(\sum_{i} (xx_{i} - x_{i}x) \otimes y_{i} = 0\) ，从而 \(xx_{i} = x_{i}x\) 对所有 \(i\) 成立（II, §3, 第 7 号, 命题 7 的推论 1）；于是必有 \(x_{i} \in \mathbb{C}'\) 对所有 \(i\) ，从而 \(z \in \mathbb{C}' \otimes_{\mathbb{K}} \mathbb{F}\) ；类似地可证 \(z \in \mathbb{E} \otimes_{\mathbb{K}} \mathbb{D}'\) ，命题得证。

推论. 若 \(\mathbf{Z}\) 和 \(\mathbf{Z}'\) 分别是 \(\mathbf{E}\) 和 \(\mathbf{F}\) 的中心，则 \(\mathbf{E} \otimes_{\mathbb{K}} \mathbf{F}\) 的中心是 \(\mathbf{Z} \otimes_{\mathbb{K}} \mathbf{Z}'\) .

设 E 和 F 是交换域 K 上代数 G 的两个子代数；假设 E 的每个元素与 F 的每个元素交换。则典范单射 \(i: E \to G\) , \(j: F \to G\) 定义一个典范同态 \(h = i \otimes j: E \otimes_{K} F \to G\) （第 2 号, 命题 5）使得

\[
(i \otimes j) (x \otimes y) = x y \quad \text { for } x \in \mathrm{E}, y \in \mathrm{F}.
\]

定义 4. 给定交换域 K 上的代数 G，若 G 的两个子代数 E 和 F 满足以下条件，则称它们在 K 上线性不相交：

\begin{enumerate}
\def\labelenumi{(\arabic{enumi})}
\item
  \(\mathbf{E}\) 的每个元素与 \(\mathbf{F}\) 的每个元素交换 ;
\item
  从 \(\mathbf{E} \otimes_{\mathbf{K}} \mathbf{F}\) 到 \(\mathbf{G}\) 的典范同态是单射。
\end{enumerate}

命题 7. 设 G 是交换域 K 上的一个代数，E, F 是 G 的两个子代数，使得 E 的每个元素与 F 的每个元素可交换。要使 E 和 F 在 K 上线性不相交，当且仅当存在 E 在 K 上的一组基，该基是 G 关于右 F-模结构的自由子集。当此条件成立时：

\begin{enumerate}
\def\labelenumi{(\roman{enumi})}
\item
  典范同态 \(h: \mathbf{E} \otimes_{\mathbb{K}} \mathbf{F} \to \mathbf{G}\) 是 \(\mathbf{E} \otimes_{\mathbb{K}} \mathbf{F}\) 到 \(\mathbf{G}\) 的由 \(\mathbf{E} \cup \mathbf{F}\) 生成的子代数上的同构 ;
\item
  \(\mathbf{E} \cap \mathbf{F} = \mathbf{K}\) ;
\item
  E（相应地，F）在 K 上的每个自由子集都是 G 的自由子集（就 G 的右或左 F-模（相应地，E-模）结构而言）。
\end{enumerate}

命题中的条件显然是必要的，因为 E 在 K 上的每一组基都是 \(E \otimes_{K} F\) 的（就其右 F-模结构而言的）基（II, §3, 第 7 号, 命题 7 的推论 1）。为看出条件是充分的，注意 h 下 \(E \otimes_{K} F\) 的像 H 是 G 中形如和 \(\sum_{i} x_{i} y_{i} = \sum_{i} y_{i} x_{i}\) 的集合，其中 \(x_{i} \in E\) 且 \(y_{i} \in F\) ；若 \((a_{\lambda})\) 是 E 在 K 上的一组基，则 H 也是由 \((a_{\lambda})\) 生成的（右或左）F-模 G 的子模。于是命题中的条件意味着存在 E 的基 \((a_{\lambda})\) ，它同时是 F-模 H 的基；由此推出 h 是单射。断言 (iii) 由 E 的每个自由子集都含于 E 的一组基中这一事实得出（II, §7, 第 1 号, 定理 2）。

推论 1. 为使从 \(\mathbf{E} \otimes_{\mathbf{K}} \mathbf{F}\) 到 \(\mathbf{G}\) 的典范同态是双射，必要且充分条件是：存在 \(\mathbf{E}\) 在 \(\mathbf{K}\) 上的一组基，它同时是（右或左）F-模 \(\mathbf{G}\) 的一组基 .

推论 2. 设 E, F 为 G 的两个子代数，在 K 上具有有限秩，且 E 的每个元素与 F 的每个元素可交换。要使 E 和 F 在 K 上线性不相交，当且仅当由 E ∪ F 生成的 G 的子代数 H 满足：

\[
[ \mathrm{H}: \mathrm{K} ] = [ \mathrm{E}: \mathrm{K} ]. [ \mathrm{F}: \mathrm{K} ].\tag{12}
\]

这就是说，满射典范同态 \(h: \mathbf{E} \otimes_{\mathbf{K}} \mathbf{F} \to \mathbf{H}\) 在 K 上的秩等于 \(\mathbf{E} \otimes_{\mathbf{K}} \mathbf{F}\) 在 \(\mathbf{K}\) 上的秩，这等价于说这个同态是双射（II, §7, 第 4 号, 命题 9）。

